\documentclass[11pt]{amsart}
\usepackage[T1]{fontenc}
\usepackage[utf8]{inputenc}
\usepackage{amsmath,amssymb,amsthm,mathtools}
\usepackage{microtype}
\usepackage[colorlinks=true,linkcolor=blue,citecolor=blue,urlcolor=blue]{hyperref}

\newtheorem{theorem}{Theorem}
\newtheorem{lemma}[theorem]{Lemma}

\theoremstyle{remark}
\newtheorem{remark}[theorem]{Remark}

\newcommand{\Cyc}[1]{C_{#1}}
\newcommand{\Z}{\mathbb{Z}}
\newcommand{\ord}{\operatorname{ord}}


\begin{document}

\title[No Cyclic $(8835,4417,2208)$ Difference Set]{On the Nonexistence of a Cyclic $(8835,4417,2208)$-Difference Set\thanks{AutoMath assisted in generating the solution and preparing this draft.}}
\date{April 15, 2026}

\begin{abstract}
We prove nonexistence of a cyclic difference set with parameters $(8835,4417,2208)$.
The argument uses the multiplier $47$, forced residue profiles modulo $3$, $5$, and $15$, and a final orbit-count contradiction on the block of elements congruent to $5$ or $10$ modulo $15$.
This removes one of the named small open cyclic Hadamard cases recorded by Baumert and Gordon.
\end{abstract}

\maketitle

\section{Introduction}
Baumert and Gordon list the parameter set $(8835,4417,2208)$ among the open cyclic Hadamard cases up to $10000$ \cite{BG2004}, and Gordon's La Jolla status slides continue to display $8835$ among the small unresolved cyclic Hadamard parameters \cite{Gordon2019}.
We show that the case is impossible.

\begin{theorem}
There is no cyclic $(8835,4417,2208)$-difference set.
\end{theorem}

\begin{proof}
Assume for contradiction that $D\subseteq G=\Cyc{8835}$ is a $(8835,4417,2208)$-difference set.
Set
\[
n=k-\lambda=4417-2208=2209=47^2.
\]
By the standard prime-power multiplier theorem, $47$ is a multiplier.
Since
\[
\gcd(47-1,8835)=\gcd(46,8835)=1,
\]
we may translate $D$ and assume
\[
47D=D.
\]
Because $47\equiv 2\pmod{m}$ for $m\in\{3,5,15\}$, the residue counts of $D$ modulo each of these moduli are constant on the multiplication-by-$2$ orbits.

\medskip
\noindent\textbf{Step 1: the profile modulo $3$.}
Let $N_i$ be the number of elements of $D$ congruent to $i$ modulo $3$.
Since multiplication by $2$ swaps the nonzero classes of $\Z/3\Z$, we have $N_1=N_2$.
For a nontrivial character of $\Z/3\Z$ with value $\omega$ on $1$, the lifted character sum is
\[
N_0+N_1\omega+N_2\omega^2 = N_0-N_1.
\]
Every nonprincipal character value has modulus $47$, so $|N_0-N_1|=47$.
Together with $N_0+2N_1=4417$, this forces
\[
(N_0,N_1,N_2)=(1441,1488,1488).
\]

\medskip
\noindent\textbf{Step 2: the profile modulo $5$.}
Let $M_j$ be the number of elements of $D$ congruent to $j$ modulo $5$.
Multiplication by $2$ acts transitively on the four nonzero residue classes, so
\[
M_1=M_2=M_3=M_4.
\]
For a nontrivial order-$5$ character with primitive value $\zeta_5$, the lifted character sum is
\[
M_0+M_1(\zeta_5+\zeta_5^2+\zeta_5^3+\zeta_5^4)=M_0-M_1.
\]
Again the modulus is $47$, so $|M_0-M_1|=47$.
Since $M_0+4M_1=4417$, we obtain
\[
(M_0,M_1,M_2,M_3,M_4)=(921,874,874,874,874).
\]

\medskip
\noindent\textbf{Step 3: the profile modulo $15$.}
Multiplication by $2$ on $\Z/15\Z$ has orbits
\[
\{0\},\quad \{5,10\},\quad \{3,6,9,12\},\quad \{1,2,4,8\},\quad \{7,11,13,14\}.
\]
Let the common counts on these five orbits be
\[
a,\ b,\ c,\ d,\ e
\]
respectively.
Comparing with the profiles modulo $3$ and $5$ found above yields
\begin{align}
a+4c &= 1441, \label{eq:mod3} \\
a+2b &= 921, \label{eq:mod5a} \\
c+d+e &= 874, \label{eq:mod5b} \\
b+2d+2e &= 1488. \label{eq:mod3b}
\end{align}

Now choose a primitive character of $\Z/15\Z$ with value $\zeta$ on $1$ and set
\[
U=\zeta+\zeta^2+\zeta^4+\zeta^8.
\]
Then the sum on the complementary unit orbit is $\overline{U}$, and one has
\[
U+\overline{U}=1,\qquad U=\frac{1\pm i\sqrt{15}}{2}.
\]
Therefore the lifted character sum is
\[
A_{15}=a-b-c+dU+e\overline{U},
\]
and $|A_{15}|=47$.
Using \eqref{eq:mod3}--\eqref{eq:mod3b}, we obtain
\[
a=1441-4c,\qquad b=2c-260,\qquad d+e=874-c,
\]
hence
\[
A_{15}=\frac{4276-15c}{2}+\frac{(d-e)i\sqrt{15}}{2}.
\]
The equation $|A_{15}|^2=47^2$ becomes
\[
(4276-15c)^2+15(d-e)^2=8836.
\]
Since $|4276-15c|\le 94$ and $4276-15c\equiv 1\pmod{15}$, there are only finitely many candidates for the first term.
A direct check shows that the unique possibility is
\[
4276-15c=-14,\qquad d-e=\pm 24.
\]
Thus
\[
c=286,\qquad b=312,\qquad a=297,\qquad \{d,e\}=\{306,282\}.
\]
In particular, the block
\[
B=\{x\in \Z/8835\Z: x\equiv 5 \text{ or } 10 \pmod{15}\}
\]
contains exactly
\[
2b=624
\]
elements of $D$.

\medskip
\noindent\textbf{Step 4: the final orbit obstruction.}
Because $47D=D$, the set $D$ is a union of full $47$-orbits in $\Cyc{8835}$.
Every element of $B$ has additive order one of
\[
3,\ 57,\ 93,\ 1767,
\]
so the possible $47$-orbit sizes inside $B$ are
\begin{align*}
\ord_3(47)&=2, & \ord_{57}(47)&=18,\\
\ord_{93}(47)&=10, & \ord_{1767}(47)&=90.
\end{align*}
The numbers of such orbits available in the whole group are
\[
\frac{\varphi(3)}{2}=1,\qquad
\frac{\varphi(57)}{18}=2,\qquad
\frac{\varphi(93)}{10}=6,\qquad
\frac{\varphi(1767)}{90}=12.
\]
Hence the total $624$ points of $D\cap B$ would have to satisfy
\[
624 = 2\alpha + 18\beta + 10\gamma + 90\delta
\]
with
\[
0\le \alpha\le 1,\qquad 0\le \beta\le 2,\qquad 0\le \gamma\le 6,\qquad 0\le \delta\le 12.
\]
After dividing by $2$,
\[
312=\alpha+9\beta+5\gamma+45\delta.
\]
If $\delta\le 4$, then the right-hand side is at most
\[
1+18+30+180=229,
\]
so $\delta\ge 5$.
If $\delta=5$, the remainder is
\[
87=\alpha+9\beta+5\gamma,
\]
but the right-hand side is at most $49$.
If $\delta=6$, the remainder is
\[
42=\alpha+9\beta+5\gamma,
\]
which is impossible for $\alpha\in\{0,1\}$, $0\le \beta\le 2$, and $0\le \gamma\le 6$.
If $\delta\ge 7$, then $45\delta>312$.
This contradiction proves that no cyclic $(8835,4417,2208)$-difference set exists.
\end{proof}

\begin{remark}
The proof isolates a rigid mod-$15$ profile and then shows that the required $\{5,10\}\pmod{15}$ block cannot be assembled from full $47$-multiplier orbits.
No search-heavy construction check is needed.
\end{remark}

\begin{thebibliography}{99}

\bibitem{BG2004}
Leonard D.~Baumert and Daniel M.~Gordon,
\emph{On the existence of cyclic difference sets with small parameters},
Fields Inst. Commun. \textbf{41} (2004), 29--40.

\bibitem{Gordon2019}
Daniel M.~Gordon,
\emph{The La Jolla Difference Set Repository},
ArasuFest talk slides, August 3, 2019.

\end{thebibliography}

\end{document}
